% Discussion and limitations (P8.7)
\section{Discussion and limitations}\label{sec:discussion}

\paragraph{Learned scores beat the expert overlay, but not by much against the simplest rule.}
The combined score improves on distance to built-up land by 0.034 AUC and on the classic AHP
overlay by 0.161. The first gap is modest because growth in Bengaluru South is overwhelmingly
edge expansion: 92\,\% of the grown cells touch land that was already built in 2018, and a single
distance already ranks such cells well. The value of the learned scores lies in the ranking
\emph{among} cells next to built-up land, where they add the composition of the neighbourhood
(how much built-up, cropland and scrub lies within 250 and \SI{500}{\m}) and terrain. Our AUC
is close to the 0.82 reported for a random forest tested on later change in Rasht
\citep{ahmadlou2016}, and far below the near-perfect values of studies that validate on a
random split of a single date \citep{li2024}: the difference comes from how the test is set
up, not from a better model.

\paragraph{Why the expert overlay does worse than one of its own criteria.} Only one of the five
MCDA criteria carries real signal on its own. The pairwise comparisons are internally
consistent (CR 0.003), but consistency says nothing about relevance: they put 70\,\% of the
weight on criteria that are close to random for this question. This is the subjectivity
\citet{malczewski2004} and \citet{mosadeghi2015} warn about, made measurable by a temporal test.

\paragraph{Similarity without labels works.} The building similarity never sees a growth label,
yet it predicts growth almost as well as the random forest trained on recent change (0.773 vs
0.779), and their combination is better than either. It also explains itself: each match
comes with the features where A and B agree and where they differ, which a planner can check.

\paragraph{Outlying growth is not predicted.} No score does better than chance for growth away
from existing built-up land, as \citet{wang2021} found for diffusional growth. Such growth
follows decisions the data do not record: land deals, new industrial estates, approvals.

\paragraph{Limitations.}
\begin{itemize}
  \item \textbf{Land-cover data.} ESRI and WorldCover disagree on built-up land (62\,\% vs
        33\,\% of the area): ESRI generalises urban areas, and it calls part of the
        Bannerghatta forest rangeland \citep{venter2022}. Growth labels inherit its errors; the
        two-year rule removes flicker but not systematic bias.
  \item \textbf{OpenStreetMap.} Road mapping more than doubled between the snapshots, mostly
        through mapping effort, so each run uses its own snapshot. OSM buildings were sparse in
        the south in 2018 \citep{herfort2023} and almost none are typed, so similarity per
        building type was not possible.
  \item \textbf{Resolution.} The 100\,m cells mix land uses; half of the top-class patches are
        single cells inside built-up areas. Cell size changes the growth label itself, so
        results at different resolutions are not directly comparable.
  \item \textbf{Past growth is not good planning.} Every score learns where building
        \emph{has} happened. Water, steep slopes and the national park are excluded, but
        zoning, flood risk, infrastructure capacity and ownership are not in the data. The
        suitability map shows where growth is \emph{likely}, as an input to planning, not a
        recommendation.
\end{itemize}
